\SourceLeaf{384}%
\noindent gibt, wenn man resp. mit
\[
\alpha,\qquad a,\qquad mmA
\]
multiplicirt und addirt:
\begin{equation}\tag*{[15]}
\left\{\begin{aligned}
a'&=\frac{am'-\alpha m''}{m}-A(aA'+bB'+cC')\\
b'&=\frac{bm'-\beta m''}{m}-B(aA'+bB'+cC')\\
c'&=\frac{cm'-\gamma m''}{m}-C(aA'+bB'+cC').
\end{aligned}\right.
\end{equation}
[Aus 6 und 8 folgt nemlich:
\[
\begin{aligned}
mmA&=b\gamma-c\beta\\
mmB&=c\alpha-a\gamma\\
mmC&=a\beta-b\alpha,
\end{aligned}
\]
und es gelten ferner die Identitäten:
\[
\begin{aligned}
m^4&=(b\gamma-c\beta)^2+mm(\alpha\alpha+aa)\\
0&=(b\gamma-c\beta)(c\alpha-a\gamma)+mm(\alpha\beta+ab)\\
0&=(b\gamma-c\beta)(a\beta-b\alpha)+mm(\alpha\gamma+ac).
\end{aligned}
\]
\indent Vermöge der Gleichungen 15 erhält man ferner aus den Gleichungen 10:]
\begin{equation}\tag*{[16]}
\left\{\begin{aligned}
a''&=\frac{\alpha m'+am''}{m}+cB'-bC'\\
b''&=\frac{\beta m'+bm''}{m}+aC'-cA'\\
c''&=\frac{\gamma m'+cm''}{m}+bA'-aB',
\end{aligned}\right.
\end{equation}
[und aus diesen Gleichungen folgt:]
\begin{equation}\tag*{[17]}
a''A'+b''B'+c''C'=-\frac{\theta m'}{m}+\frac{m''}{m}(aA'+bB'+cC'),
\end{equation}
[wo $\theta$ den Ausdruck
\[
-(\alpha A'+\beta B'+\gamma C')=-(aA''+bB''+cC'')
\]
bezeichnet. In entsprechender Weise folgt aus den Gleichungen 15:]
\begin{equation}\tag*{[18]}
a'A''+b'B''+c'C''=-\frac{\theta m'}{m}-\frac{m''}{m}(\alpha A''+\beta B''+\gamma C'').
\end{equation}
